\section{Results}\label{sec:results}

\subsection{Can the next month's diesel cost be predicted?}

\paragraph{The premium (regression): no.} No model beats the futures price as a forecast of next month's
cost (Table~\ref{tab:reg}). The best out-of-sample $R^2$ against the futures benchmark is 0.5\% (historical
mean) and none of the Diebold--Mariano tests rejects equal accuracy (all $p\ge0.11$); gradient boosting is
5\% \emph{worse} than the futures. Tuning tells the same story: in every training window the cross-validation
chose the strongest ridge penalty in the grid ($\alpha=10^4$), i.e.\ it shrank the model to a constant. This
is what market efficiency predicts \citep{fama1987,alquist2013}, and it means FuelGuard cannot promise
to buy cheaper on average.

\begin{table}[H]
\centering\small
\fittable{\tabdir/regression_oos.tex}
\caption{Regression of the premium, out of sample 2012--2026 (3\,725 days). $R^2_{OS}$ and the DM test
are relative to the futures price (zero premium); DM uses Newey--West errors with 21 lags.}
\label{tab:reg}
\end{table}

\paragraph{Spikes (classification): yes, modestly, and mostly through volatility.} Table~\ref{tab:clf}
shows that all models rank spike days better than chance (AUC 0.59--0.69, 90\% block-bootstrap intervals
all above 0.5, Appendix~\ref{app:additional}). Only two models produce \emph{probabilities} better than the
historical frequency: the random forest (Brier skill +4.1\%, interval [+1.5\%, +6.6\%]) and a one-feature
logistic regression on diesel implied volatility (+5.1\%, [+1.5\%, +8.7\%]). The full logistic and boosting
models have no probabilistic skill: their extra features add variance that the penalty removes only by
pushing predictions back to the base rate (the logistic chose the strongest penalty, $C=0.001$, in 14 of 15
years). This is the bias--variance trade-off of Appendix~\ref{app:theory}, A.1, at work: the predictable part
of the target is small, so a richer model mostly fits noise.

\begin{table}[H]
\centering\small
\fittable{\tabdir/classification_oos.tex}
\caption{Spike classification, out of sample 2012--2026 (3\,725 days, spike rate 19.4\%). Brier skill is
relative to climatology (training base rate). Precision, recall and share flagged refer to the ``elevated
risk'' flag $\hat\pi_t>\bar\pi$. Last row: days ranked by implied volatility, with no model.}
\label{tab:clf}
\end{table}

\paragraph{A disclosure.} The one-feature model was \emph{not} in our original plan. After the first run,
we found that ranking days by implied volatility alone, without any model, gives an AUC of 0.71, above every
57-feature model. We then added the parsimonious logistic to the walk-forward, with the same tuning and
leakage controls. Its result is therefore partly a finding \emph{about} our first models rather than an
independent test, but the choice is economically natural (option-implied volatility is the market's own
forecast of the size of future moves) and it was not tuned on the test years.

\begin{figure}[t]
\centering
\begin{minipage}{0.64\textwidth}\includegraphics[width=\linewidth]{fig3_auc_by_year.pdf}\end{minipage}\hfill
\begin{minipage}{0.33\textwidth}\includegraphics[width=\linewidth]{fig4_calibration.pdf}\end{minipage}
\caption{Left: AUC by test year. Right: reliability diagram (predicted vs observed spike frequency by
octile of predicted probability, pooled 2012--2026).}
\label{fig:auc}
\end{figure}

\paragraph{Stability.} Pooled numbers hide a lot of variation (Figure~\ref{fig:auc}, left). Models do well
in 2013, 2016 and 2019--2020 and badly in 2023 and 2026, when most AUCs fall to 0.2--0.4. In both years,
spike days had \emph{lower} implied volatility than the other days of the year (in 2026, 51\% against
66\%): the shocks started from relatively calm markets, and in 2026 volatility rose only after the
spike had begun, so a volatility-based signal was late. Part of the pooled AUC therefore comes from separating turbulent \emph{years} from calm ones rather
than predicting the exact month. The reliability diagram (right) shows that probabilities are roughly
calibrated but compressed: predictions above 40\% occur on fewer than 4\% of days, even in years when half of the days are spikes.

\paragraph{What drives the predictions.} For the random forest, permutation importance
(Figure~\ref{fig:imp}) is dominated by recent momentum (21-day diesel return), the price level relative to
its one-year average and volatility measures. Individual importances are small (below one AUC point) and
inflation expectations have \emph{negative} importance: shuffling them improves the test AUC, a direct sign
of overfitting. The economic reading is coherent: spikes are more likely when prices are already rising,
high relative to their history and volatile.

\begin{figure}[t]
\centering\includegraphics[width=0.92\textwidth]{fig5_importance.pdf}
\caption{Permutation importance of the random forest on the test years (mean over 2012--2026; drop in AUC
when a feature is shuffled).}
\label{fig:imp}
\end{figure}

\subsection{Does it help the haulier?}

Table~\ref{tab:bt} and Figure~\ref{fig:tradeoff} evaluate 178 monthly decisions (January 2012 -- October
2026) with a 1\% hedging margin. Never hedging costs 0.88\% more than the lockable price on average, but has a
CVaR of 16.7\%: in the worst 5\% of months the firm pays about 2\,800~EUR more than it could have locked
(26\% in the worst month). Always locking removes this risk at a cost of 1\% every month.

\begin{table}[H]
\centering\fittable{\tabdir/backtest_summary.tex}
\caption{Hedging backtest, 178 monthly decisions, 2012--2026, hedge cost 1\%. ``Mean'', ``St.dev'',
``CVaR95'' and ``Worst'' refer to the cost overrun $c_t/P_t-1$. EUR figures are for a firm using
25\,000 litres a month. ML rules: $h=1$ if $\hat\pi>2\bar\pi$, $h=\tfrac12$ if $\hat\pi>\bar\pi$, else 0.}
\label{tab:bt}
\end{table}

\begin{figure}[t]
\centering\includegraphics[width=0.85\textwidth]{fig6_tradeoff.pdf}
\caption{Average cost against tail risk of each hedging rule (lower-left is better).}
\label{fig:tradeoff}
\end{figure}

\paragraph{ML rules versus the market and naive rules.} All classification-based rules roughly halve the
tail risk of never hedging (CVaR 9.4--11.3\% instead of 16.7\%) while hedging only 21--26\% of months, and
they are cheaper on average than the market and than a constant 50\% hedge. The regression rule adds
nothing, consistently with Table~\ref{tab:reg}.

\paragraph{But a simple volatility rule does better.} Locking whenever implied volatility is above its
historical median (known at the start of each year) dominates every ML rule: lower average cost (0.49\%),
lower CVaR (6.8\%), only 8 months with an overrun above 5\%, and about 900~EUR instead of 2\,800~EUR of extra
cost in a bad month for our reference firm. It hedges more often (52\% of months) than our pre-specified ML
thresholds, which turned out too conservative. Figure~\ref{fig:budget} removes this difference by comparing
signals at the \emph{same} hedging budget: the random forest and implied volatility select bad months about
equally well, both better than spreading the same budget evenly for budgets up to 70\% of months, while momentum is worse than spreading it evenly.
In other words, \textbf{the complex model learned a noisy version of the volatility signal}, and the gain of
FuelGuard comes from hedging more when risk is high. The decision threshold matters as much as the model.
These conclusions hold with hedging margins of 0\% and 2\% (Appendix~\ref{app:additional}).

\begin{figure}[t]
\centering\includegraphics[width=0.85\textwidth]{fig8_budget.pdf}
\caption{Diagnostic (not a tradable rule: ranks use the whole test period). CVaR of the overrun when the
haulier hedges the share $q$ of months with the highest signal. ``Constant ratio'' hedges a share $q$ of
every month.}
\label{fig:budget}
\end{figure}
